\documentclass[12pt,a4paper,openany,oneside]{book}

% --- Paquetes ---
\usepackage[utf8]{inputenc}
\usepackage[spanish, es-noshorthands]{babel}
\usepackage{amsmath, amsfonts, amssymb}
\usepackage{graphicx}
\usepackage{geometry}
\usepackage{hyperref}
\usepackage{booktabs}
\usepackage{fancyhdr}
\usepackage{listings}
\usepackage{xcolor}
\usepackage{tikz}
\usepackage{pgfplots}
\usepackage{colortbl}
\usepackage{multirow}
\usepackage{hhline}
\usetikzlibrary{arrows.meta, positioning, shapes.geometric, calc, shapes, chains, scopes}
\pgfplotsset{compat=1.18}

\geometry{margin=2.5cm}

% --- Colores y estilos para código Python ---
\definecolor{codegreen}{rgb}{0,0.6,0}
\definecolor{codegray}{rgb}{0.5,0.5,0.5}
\definecolor{codepurple}{rgb}{0.58,0,0.82}
\definecolor{backcolour}{rgb}{0.95,0.95,0.92}

\lstset{
    language=Python,
    backgroundcolor=\color{backcolour},   
    commentstyle=\color{codegreen},
    keywordstyle=\color{blue},
    numberstyle=\tiny\color{codegray},
    stringstyle=\color{codepurple},
    basicstyle=\ttfamily\small,
    breaklines=true,
    frame=single,
    showstringspaces=false
}

% --- Estilo de cabecera y pie ---
\pagestyle{fancy}
\fancyhf{}
\renewcommand{\headrulewidth}{0.4pt}
\renewcommand{\footrulewidth}{0.4pt}
\fancyhead[L]{\small Carlos César Sánchez Coronel}
\fancyhead[R]{\small Sesión 11: Tareas Avanzadas de NLP}
\fancyfoot[L]{\small Carlos César Sánchez Coronel}
\fancyfoot[C]{\thepage}

% --- Portada ---
\title{
    \textbf{Sesión 11} \\ 
    \Huge \textbf{TAREAS AVANZADAS DE NLP} \\
    \Large NER, Question Answering y Summarization
}
\author{
    \small Por \\ \vspace{0.2cm}
    \Large \textsc{Carlos César Sánchez Coronel}
}
\date{2026}

\begin{document}

\maketitle
\tableofcontents
\addcontentsline{toc}{chapter}{Índice general}

% ------------------------------------------------------------------
% CAPÍTULO 11: SESIÓN 11 - TAREAS AVANZADAS DE NLP
% ------------------------------------------------------------------
\chapter{Tareas Avanzadas de NLP: NER, Question Answering y Summarization}

Los modelos pre-entrenados como BERT y GPT no son fines en sí mismos, sino herramientas poderosas para resolver tareas concretas de NLP. Esta sesión se centra en tres tareas fundamentales en la industria: el reconocimiento de entidades nombradas (NER), la respuesta automática a preguntas (Question Answering) y el resumen automático (Summarization). Se exploran sus definiciones, enfoques (extractivo vs. abstractivo), métricas de evaluación y aplicaciones prácticas, con ejemplos detallados y código de implementación.

\section{Logro de la sesión}
Aplicar modelos pre-entrenados a tareas específicas de NLP como Named Entity Recognition (NER), Question Answering y Summarization, comprendiendo las diferencias entre enfoques extractivos y abstractivos, y utilizando las métricas de evaluación adecuadas.

\section{Named Entity Recognition (NER)}

\subsection{Definición y objetivos}
El reconocimiento de entidades nombradas es la tarea de identificar y clasificar en un texto las menciones de entidades con nombre, como personas, organizaciones, lugares, fechas, cantidades, etc. Formalmente, dada una secuencia de tokens \(w_1, w_2, ..., w_n\), se desea asignar a cada token una etiqueta de entidad (o "O" para fuera de entidad).

\subsection{Sistemas de etiquetado comunes}
\begin{itemize}
    \item \textbf{BIO (Beginning, Inside, Outside)}: 
    \begin{itemize}
        \item B-PER: inicio de persona
        \item I-PER: dentro de persona
        \item B-LOC: inicio de lugar
        \item I-LOC: dentro de lugar
        \item O: fuera de entidad
    \end{itemize}
    \item \textbf{BILOU}: variante que añade L (último) y U (unitario) para mayor precisión.
\end{itemize}

\textbf{Ejemplo}:
Frase: "Steve Jobs fundó Apple en Cupertino."
Etiquetas BIO: Steve (B-PER) Jobs (I-PER) fundó (O) Apple (B-ORG) en (O) Cupertino (B-LOC) . (O)

\subsection{Arquitectura típica}
Los sistemas modernos de NER utilizan modelos pre-entrenados (como BERT) con una cabeza de clasificación por token. La representación de cada token se pasa por una capa lineal + softmax para predecir la etiqueta.

\begin{center}
\begin{tikzpicture}[
    token/.style={rectangle, draw, minimum width=1cm, minimum height=0.8cm},
    arrow/.style={-latex, thick}
]
\node[token] (t1) at (0,0) {Steve};
\node[token] (t2) at (2,0) {Jobs};
\node[token] (t3) at (4,0) {fundó};
\node[token] (t4) at (6,0) {Apple};
\node[token] (t5) at (8,0) {en};
\node[token] (t6) at (10,0) {Cupertino};

\node[token, fill=blue!20] (e1) at (0,2) {B-PER};
\node[token, fill=blue!20] (e2) at (2,2) {I-PER};
\node[token, fill=gray!20] (e3) at (4,2) {O};
\node[token, fill=green!20] (e4) at (6,2) {B-ORG};
\node[token, fill=gray!20] (e5) at (8,2) {O};
\node[token, fill=red!20] (e6) at (10,2) {B-LOC};

\draw[arrow] (t1) -- (e1);
\draw[arrow] (t2) -- (e2);
\draw[arrow] (t3) -- (e3);
\draw[arrow] (t4) -- (e4);
\draw[arrow] (t5) -- (e5);
\draw[arrow] (t6) -- (e6);
\end{tikzpicture}
\caption{Etiquetado BIO en NER.}
\end{center}

\subsection{Métrica de evaluación: F1-score por entidad}
La métrica estándar para NER es el F1-score, calculado a nivel de entidad (no de token). Una entidad se considera correcta si su tipo y sus límites coinciden exactamente con la referencia. Se calculan precisión, recall y F1 sobre el conjunto de entidades predichas.

\textbf{Ejemplo}:
Referencia: ["Steve Jobs" como PER, "Apple" como ORG, "Cupertino" como LOC]
Predicción: ["Steve Jobs" como PER, "Apple" como LOC, "Cupertino" como LOC]
\begin{itemize}
    \item Verdaderos positivos: 2 (Steve Jobs, Cupertino)
    \item Falsos positivos: 1 (Apple predicha como LOC)
    \item Falsos negativos: 1 (Apple como ORG no detectada)
\end{itemize}
Precisión = 2/3 = 0.667, Recall = 2/3 = 0.667, F1 = 0.667.

\subsection{Aplicaciones de NER}
\begin{itemize}
    \item \textbf{Búsqueda semántica}: identificar entidades en consultas para mejorar resultados.
    \item \textbf{Extracción de información}: poblar bases de datos con personas, lugares, etc.
    \item \textbf{Análisis de redes sociales}: detectar menciones de marcas o personas influyentes.
    \item \textbf{Sistemas de recomendación}: basados en entidades que interesan al usuario.
\end{itemize}

\section{Question Answering (QA)}

\subsection{Definición y tipos}
Question Answering es la tarea de responder automáticamente a preguntas formuladas en lenguaje natural. Existen dos enfoques principales:

\begin{itemize}
    \item \textbf{QA extractivo}: la respuesta es un fragmento de texto extraído directamente del contexto (documento o párrafo). Ejemplo: SQuAD (Stanford Question Answering Dataset).
    \item \textbf{QA generativo}: la respuesta se genera libremente, sin restricción de extracción. Ejemplo: sistemas de diálogo, preguntas abiertas.
\end{itemize}

\subsection{QA extractivo: arquitectura}
Sobre un modelo pre-entrenado (como BERT), se añaden dos cabezas lineales:
\begin{itemize}
    \item Una que predice la probabilidad de que cada token sea el \textbf{inicio} de la respuesta.
    \item Otra que predice la probabilidad de que sea el \textbf{fin} de la respuesta.
\end{itemize}

Dado un contexto \(C\) y una pregunta \(Q\), se construye la entrada como "[CLS] Q [SEP] C [SEP]". El modelo produce representaciones para cada token. Las puntuaciones de inicio y fin se calculan como:
\[
\text{start\_logits} = \mathbf{W}_s \mathbf{h}_i, \quad \text{end\_logits} = \mathbf{W}_e \mathbf{h}_i
\]
La pérdida es la suma de las entropías cruzadas de inicio y fin. En inferencia, se selecciona el span (i, j) que maximiza \(\text{start\_probs}[i] \times \text{end\_probs}[j]\) con \(i \le j\) y dentro de límites razonables.

\begin{center}
\begin{tikzpicture}[
    token/.style={rectangle, draw, minimum width=0.8cm, minimum height=0.8cm, font=\tiny},
    arrow/.style={-latex, thick}
]
\node[token] (cls) {[CLS]};
\node[token, right=0.2cm of cls] (q1) {¿Quién};
\node[token, right=0.2cm of q1] (q2) {fundó};
\node[token, right=0.2cm of q2] (q3) {Apple?};
\node[token, right=0.4cm of q3] (sep1) {[SEP]};
\node[token, right=0.4cm of sep1] (c1) {Steve};
\node[token, right=0.2cm of c1] (c2) {Jobs};
\node[token, right=0.2cm of c2] (c3) {fundó};
\node[token, right=0.2cm of c3] (c4) {Apple};
\node[token, right=0.2cm of c4] (c5) {en};
\node[token, right=0.2cm of c5] (c6) {Cupertino};
\node[token, right=0.2cm of c6] (sep2) {[SEP]};

\node[above=0.5cm of c1] (s1) {0.1};
\node[above=0.5cm of c2] (s2) {0.8};
\node[above=0.5cm of c3] (s3) {0.05};
\node[above=0.5cm of c4] (s4) {0.02};
\node[below=0.5cm of c1] (e1) {0.02};
\node[below=0.5cm of c2] (e2) {0.1};
\node[below=0.5cm of c3] (e3) {0.3};
\node[below=0.5cm of c4] (e4) {0.7};
\end{tikzpicture}
\caption{QA extractivo: se predicen probabilidades de inicio (arriba) y fin (abajo). La respuesta sería "Steve Jobs" (inicio en c1, fin en c2).}
\end{center}

\subsection{Métricas para QA}
\begin{itemize}
    \item \textbf{Exact Match (EM)}: porcentaje de respuestas que coinciden exactamente con la referencia (carácter por carácter).
    \item \textbf{F1-score}: a nivel de token, se calcula solapamiento de tokens entre predicción y referencia. Es más indulgente que EM.
\end{itemize}

\textbf{Ejemplo}:
Referencia: "Steve Jobs"
Predicción: "Steve"
EM = 0 (no coincide exactamente)
F1: tokens referencia {Steve, Jobs}, tokens predicción {Steve}. Precisión = 1/1 = 1, Recall = 1/2 = 0.5, F1 = 0.667.

\subsection{Aplicaciones de QA}
\begin{itemize}
    \item \textbf{Motores de respuesta}: buscadores que responden directamente (ej. Google).
    \item \textbf{Asistentes virtuales}: Alexa, Siri respondiendo preguntas.
    \item \textbf{Soporte al cliente}: responder preguntas frecuentes basadas en documentación.
\end{itemize}

\section{Summarization (Resumen Automático)}

\subsection{Definición}
El resumen automático consiste en generar una versión concisa de un texto largo, manteniendo la información más importante. Existen dos enfoques:

\begin{itemize}
    \item \textbf{Summarization extractivo}: selecciona oraciones o frases del texto original y las concatena. Es más simple y garantiza fluidez, pero puede ser menos cohesivo.
    \item \textbf{Summarization abstractivo}: genera nuevas oraciones que parafrasean y condensan la información. Es más flexible pero requiere modelos generativos complejos.
\end{itemize}

\subsection{Arquitectura para summarization abstractivo}
Se utilizan modelos seq2seq pre-entrenados (como BART, T5). El encoder procesa el documento y el decoder genera el resumen paso a paso, con atención cruzada sobre el encoder.

\begin{center}
\begin{tikzpicture}[
    layer/.style={rectangle, draw, minimum width=2.5cm, minimum height=0.8cm},
    arrow/.style={-latex, thick}
]
\node[layer, fill=blue!10] (doc) at (0,0) {Documento};
\node[layer, fill=green!10] (enc) at (0,1.5) {Encoder};
\node[layer, fill=green!10] (dec) at (0,3) {Decoder};
\node[layer, fill=orange!10] (summ) at (0,4.5) {Resumen};

\draw[arrow] (doc) -- (enc);
\draw[arrow] (enc) -- (dec);
\draw[arrow] (dec) -- (summ);
\end{tikzpicture}
\caption{Arquitectura típica para summarization abstractivo.}
\end{center}

\subsection{Métrica de evaluación: ROUGE}
ROUGE (Recall-Oriented Understudy for Gisting Evaluation) es la métrica estándar para summarization. Mide el solapamiento de n-gramas, secuencias de palabras y pares de palabras entre el resumen generado y las referencias humanas.

\textbf{Variantes comunes}:
\begin{itemize}
    \item \textbf{ROUGE-N}: solapamiento de n-gramas (generalmente unigramas y bigramas).
    \item \textbf{ROUGE-L}: basado en la subsecuencia común más larga (LCS). Captura fluidez de la oración.
    \item \textbf{ROUGE-S}: basado en skip-bigrams (pares de palabras con hasta un gap).
\end{itemize}

\textbf{Cálculo de ROUGE-1} (unigramas):
Dado un resumen generado \(G\) y una referencia \(R\):
\[
\text{ROUGE-1} = \frac{\text{\# unigramas comunes}}{\text{\# unigramas en } R} \quad (\text{recall})
\]
También se puede calcular precisión y F1.

\textbf{Ejemplo numérico}:
Referencia: "El gato persigue al ratón"
Generado: "El gato caza ratones"
Unigramas referencia: {El, gato, persigue, al, ratón}
Unigramas generado: {El, gato, caza, ratones}
Comunes: {El, gato} (2)
Recall = 2/5 = 0.4, Precisión = 2/4 = 0.5, F1 = 0.444.

ROUGE-2 (bigramas):
Bigramas referencia: {El gato, gato persigue, persigue al, al ratón}
Bigramas generado: {El gato, gato caza, caza ratones}
Comunes: {El gato} (1)
Recall = 1/4 = 0.25, Precisión = 1/3 = 0.333, F1 = 0.286.

\subsection{Desafíos en summarization}
\begin{itemize}
    \item \textbf{Fidelidad}: evitar alucinaciones (información no presente en el original).
    \item \textbf{Coherencia}: el resumen debe ser fluido y lógico.
    \item \textbf{Compresión}: decidir qué información es esencial.
\end{itemize}

\subsection{Aplicaciones de summarization}
\begin{itemize}
    \item \textbf{Resúmenes de noticias}: generación automática de titulares o resúmenes.
    \item \textbf{Informes médicos}: condensar historiales clínicos.
    \item \textbf{Documentos legales}: resumir contratos o sentencias.
    \item \textbf{Análisis de redes sociales}: resumir tendencias o conversaciones.
\end{itemize}

\section{Implementación con Hugging Face}

\subsection{NER con BERT}
\begin{lstlisting}[language=Python]
from transformers import AutoTokenizer, AutoModelForTokenClassification, pipeline

tokenizer = AutoTokenizer.from_pretrained("dslim/bert-base-NER")
model = AutoModelForTokenClassification.from_pretrained("dslim/bert-base-NER")

nlp = pipeline("ner", model=model, tokenizer=tokenizer)
text = "Steve Jobs founded Apple in Cupertino."
entities = nlp(text)
for entity in entities:
    print(entity)
\end{lstlisting}

\subsection{QA extractivo con BERT}
\begin{lstlisting}[language=Python]
from transformers import pipeline

qa_pipeline = pipeline("question-answering", model="distilbert-base-cased-distilled-squad")
context = "Steve Jobs founded Apple in Cupertino."
question = "Who founded Apple?"
result = qa_pipeline(question=question, context=context)
print(f"Answer: {result['answer']}, score: {result['score']}")
\end{lstlisting}

\subsection{Summarization abstractivo con BART}
\begin{lstlisting}[language=Python]
from transformers import pipeline

summarizer = pipeline("summarization", model="facebook/bart-large-cnn")
text = """
Elon Musk, CEO of Tesla and SpaceX, announced today that Tesla's new factory in Berlin 
will begin production of the Model Y next month. The factory, which cost over $5 billion, 
is expected to create 10,000 jobs. Musk also hinted at a new affordable electric vehicle 
to be unveiled next year.
"""
summary = summarizer(text, max_length=50, min_length=25)
print(summary[0]['summary_text'])
\end{lstlisting}

\section{Preparación para entrevistas}

\subsection{Preguntas típicas}
\begin{itemize}
    \item \textbf{"¿Qué métrica usarías para evaluar un sistema de NER?"}
    F1-score a nivel de entidad, considerando tipo y límites exactos. Se prefiere sobre accuracy por token porque este último puede ser engañoso cuando hay muchas etiquetas "O".
    
    \item \textbf{"Explica la diferencia entre QA extractivo y generativo."}
    Extractivo selecciona un fragmento del contexto; generativo produce una respuesta libre. El extractivo es más fiable pero limitado; el generativo es más flexible pero propenso a alucinaciones.
    
    \item \textbf{"¿Qué es ROUGE y cómo se calcula?"}
    Métrica para summarization basada en solapamiento de n-gramas (ROUGE-N), subsecuencias comunes (ROUGE-L) o skip-bigrams (ROUGE-S). Mide recall principalmente, pero también se usan F1.
    
    \item \textbf{"¿Qué problemas tiene el summarization abstractivo?"}
    Alucinaciones (generar información falsa), repeticiones, falta de coherencia y dificultad para capturar información clave.
    
    \item \textbf{"¿Cómo manejarías entidades anidadas en NER?"}
    Con etiquetado jerárquico o usando modelos que permitan múltiples etiquetas (ej. BERT con cabezas múltiples). También se pueden usar dos etapas: primero detectar límites, luego clasificar.
\end{itemize}

\subsection{Preguntas avanzadas}
\begin{itemize}
    \item \textbf{"¿Cómo funciona el mecanismo de atención cruzada en summarization?"}
    En modelos seq2seq, el decoder atiende a las representaciones del encoder para seleccionar información relevante del documento original. Esto permite generar resúmenes abstractivos.
    
    \item \textbf{"¿Qué es el problema de exposure bias en generación y cómo se mitiga?"}
    Ocurre cuando el modelo se entrena con teacher forcing (usa palabras verdaderas) pero en inferencia usa sus propias predicciones. Se mitiga con scheduled sampling o entrenamiento adversarial.
    
    \item \textbf{"¿Cómo evaluarías un sistema de QA sin respuestas en el contexto?"}
    Se pueden incluir preguntas sin respuesta (SQuAD 2.0). El modelo debe predecir "imposible" cuando no hay respuesta. La métrica combina EM/F1 para preguntas con respuesta y precisión para preguntas sin respuesta.
\end{itemize}

\section{Conexión con el resto del curso}
\begin{itemize}
    \item \textbf{Sesión 10 (BERT, GPT)}: Estos modelos pre-entrenados son la base de las tareas avanzadas.
    \item \textbf{Sesión 8 (Seq2Seq y Atención)}: Fundamental para summarization abstractivo.
    \item \textbf{Sesiones anteriores}: Tokenización, embeddings, fine-tuning.
\end{itemize}

\section{Resumen}
\begin{itemize}
    \item \textbf{NER}: identifica entidades en texto. Se evalúa con F1-score por entidad.
    \item \textbf{QA extractivo}: extrae respuestas de un contexto. Métricas: Exact Match y F1.
    \item \textbf{Summarization}: genera resúmenes (extractivos o abstractivos). Métrica estándar: ROUGE.
    \item Los modelos pre-entrenados (BERT, BART, T5) se fine-tunean para estas tareas con cabezas específicas.
    \item Hugging Face proporciona pipelines que simplifican la implementación.
\end{itemize}

\end{document}